\chapter{网与滤子}
\label{chap:nets-filters}

在度量空间中，许多拓扑问题都可以交给数列处理：闭包中的点是某个数列的
极限，紧致性可以表述为每个数列都有收敛子列，连续性也可由保持数列极限
刻画。一般拓扑空间未必有可数邻域基。此时，一个数列只能依次满足可数多
个邻域条件，不能保证它最终进入点的全部邻域。

网把指标集从 \(\N\) 换成任意有向集；滤子则直接记录“最终必然发生”的
集合。两种语言携带同一类渐近信息，但各有适合的论证：网保留逐项逼近的
形式，滤子在紧致性、乘积和一致结构中往往更简洁。本章建立两者的换译，
并用它们给出不附加可数性假设的闭包、连续性与紧致性判据。

\section{有向集和网}

\begin{definition}[有向集]
非空集合 \(I\) 上的关系 \(\preccurlyeq\) 若满足
\begin{enumerate}[label=\textup{(\roman*)}]
 \item \(i\preccurlyeq i\)；
 \item \(i\preccurlyeq j\preccurlyeq k\) 蕴含 \(i\preccurlyeq k\)；
 \item 对任意 \(i,j\in I\)，存在 \(k\in I\) 使
 \(i\preccurlyeq k\) 且 \(j\preccurlyeq k\)，
\end{enumerate}
则称 \((I,\preccurlyeq)\) 为\term{有向集}。这里不要求反对称性。
\end{definition}

第三条只要求任意有限组条件能够同时向后推进。通常的
\((\N,\le)\) 是有向集；一个集合的有限子集全体按包含关系排序也是有向
集；若 \(x\) 是拓扑空间中的点，则 \(x\) 的邻域全体按反向包含
\[
 U\preccurlyeq V\quad\Longleftrightarrow\quad V\subseteq U
\]
排序同样是有向集。最后一个例子已经表明，有向关系表示“信息变得更精细”，
不一定表示数值增大。

\begin{definition}[网]
设 \(I\) 为有向集，映射 \(x:I\to X\) 称为 \(X\) 中的一个\term{网}，
记作 \((x_i)_{i\in I}\)。若性质 \(P(i)\) 对某个 \(i_0\in I\) 之后
的全部指标成立，即
\[
 i\succcurlyeq i_0\quad\Longrightarrow\quad P(i),
\]
则称 \(P(i)\) \term{最终成立}；若对每个 \(i_0\) 都存在
\(i\succcurlyeq i_0\) 使 \(P(i)\) 成立，则称它\term{频繁成立}。
\end{definition}

数列是指标集为 \(\N\) 的网。更能体现一般性的例子来自有限逼近。设
\(S\) 是任意集合，\(\mathcal D\) 为 \(S\) 的有限子集全体，按包含排序。
若对每个 \(F\in\mathcal D\) 构造一个同时满足 \(F\) 中全部条件的对象
\(x_F\)，那么“对每个固定条件 \(s\in S\)，当 \(F\) 充分大时该条件
成立”就成为普通的最终性陈述。

\begin{definition}[网的收敛]
设 \(X\) 为拓扑空间。网 \((x_i)_{i\in I}\) \term{收敛}到
\(x\in X\)，若对 \(x\) 的每个邻域 \(U\)，都有 \(x_i\in U\)
最终成立，记作 \(x_i\to x\)。
\end{definition}

若把 \(x\) 的邻域全体作为指标集，并在每个邻域 \(U\) 中选取一点
\(x_U\)，则所得网自动趋于 \(x\)。为了避免在任意多邻域上预先作选择，
也可把指标写成
\[
 I=\{(U,y):U\text{ 是 }x\text{ 的邻域},\ y\in U\},
\]
并令 \(x_{(U,y)}=y\)。按邻域反向包含排序后，它同样趋于 \(x\)。

\begin{proposition}[极限的唯一性]
\label{prop:net-limits-hausdorff}
拓扑空间 \(X\) 为 Hausdorff 空间，当且仅当 \(X\) 中每个收敛网的
极限唯一。
\end{proposition}

\begin{proof}
若 \(X\) 为 Hausdorff 空间，而同一网同时趋于 \(x\ne y\)，取不交邻域
\(U\ni x\)、\(V\ni y\)。网最终在 \(U\) 中，也最终在 \(V\) 中；有向
性使这两个最终条件可以同时满足，得到某项属于 \(U\cap V\)，矛盾。

反之，若 \(x\ne y\) 不能由不交邻域分开，令指标集由三元组
\((U,V,z)\) 组成，其中 \(U\ni x\)、\(V\ni y\) 为邻域且
\(z\in U\cap V\)，按 \(U,V\) 同时反向包含排序，并取网值为 \(z\)。
这个网既趋于 \(x\) 又趋于 \(y\)，故极限不唯一。
\end{proof}

\section{网的收敛及子网}

数列的子列由严格递增指标给出。一般有向集没有一个适用的“严格递增整数
指标”，需要直接规定新指标最终越过旧指标的每个阶段。

\subsection{最终共尾映射与聚点}

有向集 \(I\) 的子集 \(D\) 称为\term{共尾子集}，若对每个 \(i\in I\)，
都存在 \(d\in D\) 使 \(i\preccurlyeq d\)。它带诱导关系也有向：先在
\(I\) 中取两个指标的共同上界，再取超过该上界的 \(D\) 中元素。
仅要求一个映射的像是共尾子集，还不能保证它把后期指标送到后期指标。

\begin{definition}[子网]
设 \((x_i)_{i\in I}\) 是网。若 \(J\) 是有向集，映射
\(\phi:J\to I\) 满足
\[
 \forall i_0\in I\ \exists j_0\in J\ \forall j\succcurlyeq j_0,
 \qquad \phi(j)\succcurlyeq i_0,
\]
则称 \(\phi\) 为\term{最终共尾映射}，并称
\((x_{\phi(j)})_{j\in J}\) 为 \((x_i)\) 的一个\term{子网}。
\end{definition}

有些文献还要求 \(\phi\) 保序。本书采用上面的最终共尾定义；它恰好保证
原网最终成立的每个性质在子网上仍最终成立。因此，收敛网的每个子网收敛
到同一极限。

例如令 \(\phi:\N\to\N\) 满足 \(\phi(2n)=n\)、\(\phi(2n+1)=0\)。
其像为整个 \(\N\)，当然共尾，但它不是最终共尾映射，因为任意晚的奇数
指标仍被送到 \(0\)。因此它不能作为上述子网定义中的指标映射。

\begin{definition}[网的聚点]
点 \(x\in X\) 称为网 \((x_i)\) 的\term{聚点}，若对 \(x\) 的每个
邻域 \(U\) 和每个 \(i_0\in I\)，都存在 \(i\succcurlyeq i_0\) 使
\(x_i\in U\)。也就是说，\(x_i\in U\) 对每个邻域 \(U\) 都频繁成立。
\end{definition}

\begin{theorem}[聚点与收敛子网]
\label{thm:net-cluster-subnet}
点 \(x\) 是网 \((x_i)_{i\in I}\) 的聚点，当且仅当该网有一个收敛
到 \(x\) 的子网。
\end{theorem}

\begin{proof}
若子网 \(x_{\phi(j)}\to x\)，给定邻域 \(U\) 和 \(i_0\)，子网最终进入
\(U\)，而 \(\phi\) 最终越过 \(i_0\)。把两个最终条件合并，便得到某个
\(i=\phi(j)\succcurlyeq i_0\) 且 \(x_i\in U\)。

反过来，设 \(x\) 是聚点。令
\[
 J=\{(i,U):i\in I, U\text{ 是 }x\text{ 的邻域},\ x_i\in U\}.
\]
规定
\[
 (i,U)\preccurlyeq(j,V)
 \quad\Longleftrightarrow\quad i\preccurlyeq j\text{ 且 }V\subseteq U.
\]
给定 \((i,U),(j,V)\)，先取 \(k_0\succcurlyeq i,j\)，再由聚点性质取
\(k\succcurlyeq k_0\) 使 \(x_k\in U\cap V\)。于是
\((k,U\cap V)\) 是共同上界，故 \(J\) 有向。令
\(\phi(i,U)=i\)。聚点性质保证 \(J\) 非空；给定 \(i_0\)，取某个
\(i\succcurlyeq i_0\) 和邻域 \(X\) 得到 \((i,X)\in J\)，其后的
指标都被 \(\phi\) 送到 \(i_0\) 之后，所以 \(\phi\) 最终共尾。
最后，给定邻域 \(U_0\)，从某个 \((i,U_0)\in J\) 开始，后继
\((j,V)\) 满足 \(x_j\in V\subseteq U_0\)。所得子网收敛到 \(x\)。
\end{proof}

网之所以能够恢复一般拓扑，关键是指标集可以同时编码全部邻域。

\subsection{闭包与连续性}

\begin{theorem}[闭包的网刻画]
\label{thm:closure-by-nets}
设 \(A\subseteq X\)。则
\[
 x\in\overline A
 \quad\Longleftrightarrow\quad
 \text{存在 }A\text{ 中的网收敛到 }x.
\]
因而，\(A\) 闭当且仅当 \(A\) 中每个收敛网的极限仍在 \(A\) 中。
\end{theorem}

\begin{proof}
若 \(x\in\overline A\)，每个邻域 \(U\) 都与 \(A\) 相交。取指标集
\[
 I=\{(U,a):U\ni x\text{ 为邻域},\ a\in A\cap U\},
\]
按邻域反向包含排序，并令 \(x_{(U,a)}=a\)。这给出 \(A\) 中趋于
\(x\) 的网。反向结论直接来自闭包 \(\overline A\) 是含 \(A\) 的闭集：
若网始终在 \(A\) 中并趋于 \(x\)，则每个 \(x\) 的邻域都与 \(A\)
相交，故 \(x\in\overline A\)。闭集刻画由
\(A=\overline A\) 得到。
\end{proof}

\begin{theorem}[连续性的网判据]
\label{thm:continuity-by-nets}
映射 \(f:X\to Y\) 在 \(x\in X\) 连续，当且仅当对每个
\(x_i\to x\)，都有 \(f(x_i)\to f(x)\)。
\end{theorem}

\begin{proof}
连续性立即推出保持网的极限。反之，若 \(f\) 在 \(x\) 不连续，则存在
\(f(x)\) 的邻域 \(V\)，使 \(f^{-1}(V)\) 不是 \(x\) 的邻域。
因此每个 \(x\) 的邻域 \(U\) 都含有一点
\(x_U\notin f^{-1}(V)\)。用成对指标 \((U,y)\) 可避免选择：令
\[
 I=\{(U,y):U\ni x, y\in U\setminus f^{-1}(V)\},
\]
按 \(U\) 反向包含排序。网 \(y_{(U,y)}=y\) 趋于 \(x\)，但其像网
没有任何一项落在 \(V\) 中，故不趋于 \(f(x)\)，与假设矛盾。
\end{proof}

\begin{example}[序列不能刻画的闭包]
取不可数集合 \(S\)，令 \(X=\{0,1\}^{S}\) 取乘积拓扑，\(A\) 为所有
只有有限多个坐标等于 \(1\) 的点。恒等于 \(1\) 的点记为 \(\boldsymbol1\)。
每个基本邻域只限制有限多个坐标，故与 \(A\) 相交，于是
\(\boldsymbol1\in\overline A\)。

然而，\(A\) 中任一数列 \((a_n)\) 的支集并集仍至多可数。可取
\(s\in S\) 不在这个并集中，于是 \(a_n(s)=0\) 对所有 \(n\) 成立，
该数列不可能趋于 \(\boldsymbol1\)。另一方面，以 \(S\) 的有限子集
\(F\) 为指标，令 \(a_F=\boldsymbol1_F\)，则 \(a_F\to\boldsymbol1\)。
\end{example}

\section{滤子和滤子基}

\begin{definition}[滤子]
集合 \(X\) 上的一个\term{滤子}是子集族 \(\mathcal F\subseteq
\mathcal P(X)\)，满足
\begin{enumerate}[label=\textup{(\roman*)}]
 \item \(\varnothing\notin\mathcal F\)，且 \(X\in\mathcal F\)；
 \item \(A,B\in\mathcal F\) 蕴含 \(A\cap B\in\mathcal F\)；
 \item \(A\in\mathcal F\) 且 \(A\subseteq B\subseteq X\) 蕴含
 \(B\in\mathcal F\)。
\end{enumerate}
若 \(\mathcal F\subseteq\mathcal G\)，称 \(\mathcal G\) 比
\(\mathcal F\) \term{更细}。
\end{definition}

滤子中的集合可理解为最终必须进入的集合。点 \(x\) 的邻域滤子
\(\mathcal N(x)\) 由所有含 \(x\) 的邻域组成；网 \((x_i)\) 的
\term{尾滤子}为
\[
 \mathcal F_x=\{A\subseteq X:\text{存在 }i_0,
 \{x_i:i\succcurlyeq i_0\}\subseteq A\}.
\]

\subsection{滤子基与网的换译}

\begin{definition}[滤子基]
非空集合族 \(\mathcal B\subseteq\mathcal P(X)\setminus\{\varnothing\}\)
称为\term{滤子基}，若对任意 \(B_1,B_2\in\mathcal B\)，存在
\(B_3\in\mathcal B\) 使 \(B_3\subseteq B_1\cap B_2\)。它生成的滤子为
\[
 \langle\mathcal B\rangle
 =\{A\subseteq X:\text{存在 }B\in\mathcal B, B\subseteq A\}.
\]
\end{definition}

网的尾部集合构成滤子基。反过来，给定滤子基 \(\mathcal B\)，令
\[
 I_{\mathcal B}=\{(B,x):B\in\mathcal B, x\in B\},
\]
并按 \(B\) 的反向包含排序，取 \(x_{(B,x)}=x\)，便得到与
\(\langle\mathcal B\rangle\) 携带同样最终信息的网。使用成对指标使这个
构造不需要为每个 \(B\) 预先选点。

更精确地，从 \((B,x)\) 开始的网值集合恰为 \(B\)。一方面，后继
\((C,y)\) 满足 \(y\in C\subseteq B\)；另一方面，每个 \(y\in B\)
都出现在后继指标 \((B,y)\) 处。因此这个网的尾滤子就是
\(\langle\mathcal B\rangle\)，不仅有相同的收敛点。

\subsection{收敛与聚集}

\begin{definition}[滤子的收敛与聚点]
滤子 \(\mathcal F\) \term{收敛}到 \(x\)，若
\(\mathcal N(x)\subseteq\mathcal F\)，记作 \(\mathcal F\to x\)。
点 \(x\) 是 \(\mathcal F\) 的\term{聚点}，若
\[
 A\cap U\ne\varnothing
 \quad(A\in\mathcal F, U\in\mathcal N(x)).
\]
等价地，\(x\in\bigcap_{A\in\mathcal F}\overline A\)。
\end{definition}

对网的尾滤子，滤子收敛等价于网收敛，滤子聚点等价于网聚点。更细的滤子
保留收敛，但未必保留聚点。例如在二点离散空间 \(X=\{a,b\}\) 中，
滤子 \(\{X\}\) 的聚点为 \(a,b\)；较细的主滤子
\(\{A\subseteq X:a\in A\}\) 只有聚点 \(a\)，因为其成员 \(\{a\}\)
与 \(b\) 的邻域 \(\{b\}\) 不交。

\subsection{超滤子}

\begin{definition}[超滤子]
在包含关系下极大的滤子称为\term{超滤子}。
\end{definition}

\begin{proposition}[超滤子的二分性质]
\label{prop:ultrafilter-dichotomy}
滤子 \(\mathcal U\) 是超滤子，当且仅当对每个 \(A\subseteq X\)，
恰有 \(A\in\mathcal U\) 与 \(X\setminus A\in\mathcal U\) 之一成立。
\end{proposition}

\begin{proof}
两者不能同时成立，否则空集属于 \(\mathcal U\)。若两者都不属于
\(\mathcal U\)，则每个 \(F\in\mathcal U\) 都与 \(A\) 相交；否则
\(F\subseteq X\setminus A\) 会推出补集属于 \(\mathcal U\)。于是
\(\{F\cap A:F\in\mathcal U\}\) 是滤子基，生成严格包含
\(\mathcal U\) 的滤子，违背极大性。反向结论也直接：若
\(\mathcal G\supsetneq\mathcal U\) 是滤子，取
\(A\in\mathcal G\setminus\mathcal U\)，则 \(X\setminus A\in\mathcal U
\subseteq\mathcal G\)，从而空集属于 \(\mathcal G\)，矛盾。
\end{proof}

\begin{theorem}[超滤子扩张原理]
\label{thm:ultrafilter-extension}
集合 \(X\) 上的每个滤子都包含在某个超滤子中。
\end{theorem}

\begin{proof}
设原滤子为 \(\mathcal F\)，令 \(P\) 为所有包含
\(\mathcal F\) 的真滤子所成的集合，按包含关系排序。它非空，因为
\(\mathcal F\in P\)。对非空链 \(\mathcal C\subseteq P\)，
令 \(\mathcal H=\bigcup_{\mathcal G\in\mathcal C}\mathcal G\)。
它含 \(X\)，不含空集，且向上封闭。若 \(A,B\in\mathcal H\)，链中
分别含 \(A,B\) 的两个滤子可比较，较大者同时含二者及其交；所以
\(A\cap B\in\mathcal H\)。因此 \(\mathcal H\) 仍为真滤子，是该链
在 \(P\) 中的上界。空链则以 \(\mathcal F\) 为上界。

由Zorn引理，\(P\) 有极大元 \(\mathcal U\)。若存在严格更细
的真滤子，它仍包含 \(\mathcal F\)，便与此极大性矛盾。故 \(\mathcal U\)
是超滤子。
\end{proof}

这里使用了附录\ref{chap:appendix-set-theory}中陈述的Zorn引理。上述证明在全书采用的
选择公理背景下成立，不应把无限链的极大元存在性当作有限选择的结果。

\section{紧致性、连续性与三种收敛语言}

\subsection{紧致性的等价判据}

\begin{theorem}[紧致性的网与滤子刻画]
\label{thm:compactness-nets-filters}
对拓扑空间 \(X\)，下列命题等价：
\begin{enumerate}[label=\textup{(\roman*)}]
 \item \(X\) 紧致；
 \item \(X\) 中每个滤子都有聚点；
 \item \(X\) 中每个超滤子都收敛；
 \item \(X\) 中每个网都有收敛子网。
\end{enumerate}
\end{theorem}

空空间也是本定理的一种情形：它紧致，却没有真滤子，也没有以非空有向集
为指标的网，因此后三条均为空真。

\begin{proof}
\((\mathrm i)\Rightarrow(\mathrm{ii})\)。对滤子 \(\mathcal F\)，闭集族
\(\{\overline A:A\in\mathcal F\}\) 具有有限交性质。事实上，滤子中的
有限交非空，并且
\[
 \overline{A_1\cap\cdots\cap A_n}
 \subseteq\overline A_1\cap\cdots\cap\overline A_n.
\]
紧致性给出
\[
 \bigcap_{A\in\mathcal F}\overline A\ne\varnothing,
\]
其中任一点都是聚点。

\((\mathrm{ii})\Rightarrow(\mathrm{iii})\)。设超滤子 \(\mathcal U\)
有聚点 \(x\)。若某邻域 \(V\ni x\) 不属于 \(\mathcal U\)，二分性质
给出 \(X\setminus V\in\mathcal U\)，这与聚点条件矛盾。因此
\(\mathcal N(x)\subseteq\mathcal U\)，即 \(\mathcal U\to x\)。

\((\mathrm{iii})\Rightarrow(\mathrm{ii})\)。把 \(\mathcal F\) 扩张为
超滤子 \(\mathcal U\)。若 \(\mathcal U\to x\)，则 \(x\) 是
\(\mathcal U\) 的聚点，更是较粗滤子 \(\mathcal F\) 的聚点。

\((\mathrm{ii})\Rightarrow(\mathrm{iv})\)。任一网的尾滤子有聚点；该点
也是网的聚点，由定理 \ref{thm:net-cluster-subnet} 得到收敛子网。

\((\mathrm{iv})\Rightarrow(\mathrm i)\)。若开覆盖 \(\mathcal O\) 没有
有限子覆盖，令
\[
 I=\{(\mathcal E,x):\mathcal E\subseteq\mathcal O\text{ 有限},
 x\notin\bigcup\mathcal E\},
\]
按有限族 \(\mathcal E\) 的包含关系定向，并取网值为第二坐标 \(x\)。
这是非空有向集：两个有限族的并仍不能覆盖 \(X\)。若某子网趋于
\(x_0\)，取 \(O_0\in\mathcal O\) 使 \(x_0\in O_0\)。共尾性保证子网
最终对应于含 \(O_0\) 的有限族，因此其值最终不在 \(O_0\) 中；收敛性
却要求它最终在 \(O_0\) 中，矛盾。故 \(\mathcal O\) 必有有限子覆盖。
\end{proof}

对映射 \(f:X\to Y\) 和滤子 \(\mathcal F\)，定义像滤子
\[
 f_*\mathcal F
 =\{B\subseteq Y:f^{-1}(B)\in\mathcal F\}.
\]
它是由集合族 \(\{f(A):A\in\mathcal F\}\) 生成的滤子。

\subsection{连续映射与可数性}

\begin{corollary}[连续性的三种表述]
映射 \(f:X\to Y\) 在 \(x\) 连续，当且仅当下列任一条件成立：
\begin{enumerate}[label=\textup{(\roman*)}]
 \item 对每个网 \(x_i\to x\)，有 \(f(x_i)\to f(x)\)；
 \item 对每个滤子 \(\mathcal F\to x\)，有
 \(f_*\mathcal F\to f(x)\)；
 \item \(f_*\mathcal N(x)\to f(x)\)。
\end{enumerate}
\end{corollary}

\begin{proof}
网的表述已由定理 \ref{thm:continuity-by-nets} 证明。若 \(f\) 连续，
则每个 \(f(x)\) 的邻域 \(V\) 都满足
\(f^{-1}(V)\in\mathcal N(x)\subseteq\mathcal F\)，故
\(V\in f_*\mathcal F\)。取 \(\mathcal F=\mathcal N(x)\) 得到第三条。
反之，第三条说明每个 \(f(x)\) 的邻域 \(V\) 都有
\(f^{-1}(V)\in\mathcal N(x)\)，正是点连续性的邻域定义。
\end{proof}

若 \(x\) 有递减的可数邻域基 \((U_n)\)，则从
\(x\in\overline A\) 可逐项取 \(a_n\in A\cap U_n\)，得到
\(a_n\to x\)。因此第一可数空间中的闭包与连续性仍可用数列刻画。
紧致性的序列刻画还需要额外条件；在度量空间中，这些条件自动满足。一般
拓扑中，网或滤子的判据不需要补充可数性假设。

\section*{习题}
\addcontentsline{toc}{section}{习题}

\subsection*{A组　网与滤子的基本运算}
\begin{enumerate}[label=\textbf{\thechapter.\arabic*.}]
 \exerciseitem{directed-finite-subsets}
 设 \(S\) 为集合。证明其有限子集全体按包含关系构成有向集，并说明当
 \(S\) 不可数时，该有向集不存在可数共尾子集。
 \exerciseitem{net-eventual-frequent-negation}
 证明：性质 \(P\) 频繁成立当且仅当“\(P\) 最终不成立”为假；并说明
 “\(P\) 不最终成立”不等于“\(P\) 最终不成立”。
 \exerciseitem{subnet-preserves-limit}
 从最终共尾映射的定义证明收敛网的每个子网收敛到同一极限。
 \exerciseitem{closure-neighborhood-net}
 不预先选取邻域中的点，用成对指标重新证明闭包的网刻画。
 \exerciseitem{tail-filter-verification}
 证明任一网的尾部集合生成一个滤子，并证明网收敛等价于其尾滤子收敛。
 \exerciseitem{filter-base-canonical-net}
 验证正文中 \(I_{\mathcal B}\) 有向，并证明其规范网与
 \(\langle\mathcal B\rangle\) 有相同的收敛点。
\end{enumerate}

\subsection*{B组　聚点、连续与紧致}
\begin{enumerate}[label=\textbf{\thechapter.\arabic*.},resume]
 \exerciseitem{net-cluster-tail-closure}
 证明 \(x\) 是网 \((x_i)\) 的聚点，当且仅当
 \[
  x\in\bigcap_{i_0\in I}\overline{\{x_i:i\succcurlyeq i_0\}}.
 \]
 \exerciseitem{ultrafilter-cluster-converges}
 设 \(\mathcal U\) 为超滤子。证明 \(x\) 是 \(\mathcal U\) 的聚点
 当且仅当 \(\mathcal U\to x\)。
 \exerciseitem{continuous-filter-image}
 证明像滤子的两种定义一致：
 \[
 f_*\mathcal F=\{B:f^{-1}(B)\in\mathcal F\}
 \]
 等于由 \(\{f(A):A\in\mathcal F\}\) 生成的滤子。
 \exerciseitem{compact-closed-subspace-filters}
 用滤子刻画证明：紧致空间的闭子空间紧致；Hausdorff 空间的紧致子集闭。
 \exerciseitem{product-coordinate-net-convergence}
 设 \(X=\prod_{s\in S}X_s\) 取乘积拓扑。证明网 \(x_i\to x\) 当且仅当
 对每个 \(s\in S\)，坐标网 \(x_i(s)\to x(s)\)。
 \exerciseitem{first-countable-sequential-closure}
 设 \(X\) 第一可数。证明 \(A\) 闭当且仅当 \(A\) 中每个收敛数列的
 极限仍在 \(A\) 中，并由此证明序列连续等价于连续。
\end{enumerate}

\subsection*{C组　不可数坐标与选择原则}
\begin{enumerate}[label=\textbf{\thechapter.\arabic*.},resume]
 \exerciseitem{finite-support-product-closure}
 把正文有限支集例子的每一步写完整：证明
 \(\boldsymbol1\in\overline A\)，构造趋于它的网，并证明不存在
 \(A\) 中趋于它的数列。
 \optionalexerciseitem{universal-subnet-ultrafilter}
 设 \((x_i)\) 为网，把其尾滤子扩张为超滤子 \(\mathcal U\)。由
 \(\mathcal U\) 的一个滤子基构造原网的子网，并证明该子网对每个
 \(A\subseteq X\) 都最终位于 \(A\) 或最终位于 \(X\setminus A\)。
 \optionalexerciseitem{compact-map-closed}
 设 \(X\) 紧致、\(Y\) Hausdorff，\(f:X\to Y\) 连续。用网证明
 \(f\) 是闭映射；若 \(f\) 还是双射，证明它是同胚。
\end{enumerate}
